% --- Title page -------------------------------------------------------------
% Appended after the preset partials, which is the only place it can win: the
% partials set \pretitle/\posttitle themselves, and whichever definition comes
% last is the one `titling' uses.
%
% This is why the render goes through pandoc here rather than through
% mcp-latex's render_markdown_to_pdf: that tool composes its header from
% common + type + layout and passes it as --include-in-header, which displaces
% a document's own `header-includes'.  There is no caller hook to append to.
\usepackage{tikz}
\usepackage{titling}
\definecolor{mcpaccent}{HTML}{5A3E8E}

% The mark: a bracket pair around a block cursor -- Emacs Lisp in two glyphs.
% Drawn rather than shipped as an asset so it scales cleanly and lives in
% version control as text.  The argument is a length and scales the whole
% glyph, so every coordinate below is a multiple of it.
\newcommand{\mcpmark}[1][22pt]{%
  \begin{tikzpicture}[x=#1,y=#1,baseline]
    \draw[mcpaccent,line width=0.085*#1,line cap=round,line join=round]
      (0.42,0)--(0.16,0)--(0.16,1)--(0.42,1);
    \draw[mcpaccent,line width=0.085*#1,line cap=round,line join=round]
      (1.30,0)--(1.56,0)--(1.56,1)--(1.30,1);
    \fill[mcpaccent] (0.66,0.24) rectangle (1.06,0.76);
  \end{tikzpicture}%
}

% A hairline across the type area, above and below the title.
\newcommand{\mcprule}{{\color{mcpaccent}\rule{\linewidth}{0.4pt}}}

\makeatletter
% The version arrives in the `date' field (render-reference.sh stamps it there)
% and the render date comes from __DOC_STAMP__ in the type partial.  Both are
% wanted, in that order, inside the block below -- so the version is captured
% at \begin{document}, once pandoc's \date{} from the metadata has run but
% before \maketitle reads it.
\AtBeginDocument{%
  \let\mcpversion\@date
  \let\mcpauthor\@author
  \date{}\author{}%
}
\pretitle{%
  \begin{center}
  \vspace*{1.5cm}
  \mcpmark[26pt]\\[2.4em]
  \mcprule\\[1.6em]
  \huge\bfseries
}
% Everything under the title is emitted here, so author and stamp appear once,
% in one block, rather than through the partial's own hooks.
\posttitle{%
  \par\vspace{0.9em}
  \mcprule
  \end{center}
  \begin{center}
    \vspace{1.6em}
    \large\itshape\@subtitle\par
  \end{center}
  \vfill
  \begin{center}
    {\scshape\large\mcpauthor}\par
    \vspace{0.7em}
    {\color{mcpaccent}\small\mcpversion\ \textperiodcentered\ \mcpl@stamp}\par
    \vspace{1.3em}
    \mcprule
  \end{center}
  \vspace*{1.5cm}
}
% The partial's subtitle and stamp blocks hang off its \posttitle, which is
% replaced above; blank them so a future partial change cannot reintroduce a
% second copy.
\renewcommand{\mcpl@subtitleblock}{}
\renewcommand{\mcpl@stampblock}{}
% titling prints the author and date through \preauthor/\postauthor and
% \predate/\postdate.  Both are already in the block above, and \author{} /
% \date{} at \begin{document} leaves these pairs with nothing to wrap -- but
% blanking them too keeps the vertical space they would still contribute from
% pushing the closing rule down the page.
\preauthor{}
\postauthor{}
\predate{}
\postdate{}
\makeatother
